\section*{CHAPTER 2. SYSTEM ARCHITECTURE AND MATHEMATICAL FORMULATION}
\addcontentsline{toc}{section}{\numberline{}CHAPTER 2. SYSTEM ARCHITECTURE AND MATHEMATICAL FORMULATION}

\setcounter{section}{2}
\setcounter{subsection}{0}
\setcounter{figure}{0}
\setcounter{table}{0}


\subsection{System Parameters, Sets, and Indices}

To ensure realistic simulations and validate the proposed mathematical models, the system parameters and input datasets are constructed based on physical configurations. The foundational global parameters and component data specifications are detailed as follows:

\begin{itemize}
    \item \textbf{Global System Settings:} The multi-microgrid network is normalized using a system base apparent power ($S_{base}$) of 1 MVA. The operational scheduling is performed over a day-ahead time horizon denoted by $\mathcal{T}$, comprising 24 intervals with a discrete time step of $\Delta t = 1$ hour.
    
    \item \textbf{Peer-to-Peer Trading Limits:} To respect the physical thermal capacities of the interconnected cables, the active power exchanged via tie-lines ($P_{tie,i,t}$) between interconnected microgrids is strictly bounded within a trading limit of $\pm 1.5$ MW (equivalent to 1.5 pu).
    
    \item \textbf{VOLL:} In severe emergency scenarios, load shedding may be unavoidable. However, to guarantee the absolute protection of critical loads, VOLL associated with critical load shedding ($P_{shed,i,t}^{\text{critical}}$) is assigned a strict penalty (e.g., 100,000 cents/MWh). This heavy penalty strictly forces the optimization engine to ensure 100\% continuous supply to critical demands.
    
    \item \textbf{Input Data Profiles and Network Attributes:} The dynamic 24-hour profiles for active ($P_{load,i,t}$) and reactive ($Q_{load,i,t}$) power demand, alongside the maximal generation capacities for renewable sources ($P_{PV,i,t}^{\text{peak}}$ and $P_{WT,i,t}^{\text{peak}}$), are derived from physical load and source datasets. Furthermore, the structural network parameters, including the resistance ($R_{ij}$) and reactance ($X_{ij}$) of all internal distribution lines and tie-lines, are directly derived from detailed line configuration data to accurately reflect the real-world network topology.
\end{itemize}

To formulate the mathematical models of the proposed multi-microgrid system, several key sets and indices are defined. The operational time horizon is discretized into a set of intervals denoted by $\mathcal{T}$. Let $\mathcal{M}$ denote the set of all microgrids in the network. Within the system, the distribution network is represented by a set of nodes $\mathcal{N}$ and a set of connecting distribution lines (branches) $\mathcal{E}$. The specific sets of nodes equipped with distributed energy resources are defined as follows: $\mathcal{N}^{PV}$ for Photovoltaic systems, $\mathcal{N}^{WT}$ for Wind Turbines, $\mathcal{N}^{DG}$ for Diesel Generators, and $\mathcal{N}^{BESS}$ for Battery Energy Storage Systems. 

Furthermore, the set of boundary nodes that possess tie-lines enabling physical Peer-to-Peer active power trading with neighboring microgrids is denoted by $\mathcal{N}^{trade}$. For the radial topology of the distribution network, $\Omega(i)$ represents the set of all child nodes directly connected to a parent node $i$, and $\mathcal{C}$ defines the set of available line configurations. Throughout the formulation, subscripts $i$ and $j$ refer to spatial node or branch indices, whereas the subscript $t$ refers to the discrete time step.

\subsection{Mathematical Modeling of Energy Components}

This section details the mathematical formulation governing the operational constraints of the various energy components deployed within the microgrid network.

\subsubsection{Renewable Energy Sources and Loads}
The integration of RES, specifically PV systems and WT, alongside power loads, constitutes the primary non-dispatchable elements of the network. The active power output from PV and WT units is strictly bounded by their available forecasted generation peaks at any given time step $t$. The active power limits are formulated as follows:
\begin{subequations}
\begin{align}
    0 \le P_{PV,i,t} &\le P_{PV,i,t}^{\text{peak}}, \quad \forall i \in \mathcal{N}^{PV}, \forall t \label{eq:pv_limit} \\
    0 \le P_{WT,i,t} &\le P_{WT,i,t}^{\text{peak}}, \quad \forall i \in \mathcal{N}^{WT}, \forall t \label{eq:wt_limit}
\end{align}
\end{subequations}
where $P_{PV,i,t}$ and $P_{WT,i,t}$ denote the scheduled active power of the PV and WT at node $i$ during time $t$, respectively. $P_{PV,i,t}^{\text{peak}}$ and $P_{WT,i,t}^{\text{peak}}$ represent their maximum available active power based on weather conditions.

For the load demand, both active ($P_{L,i,t}$) and reactive ($Q_{L,i,t}$) power are treated as fixed parameters unless load shedding is activated during emergency operations.

\subsubsection{BESS}
The BESS provides essential flexibility by charging during surplus periods and discharging during generation shortages. To prevent simultaneous charging and discharging, binary indicators $u_{ch,i,t}$ and $u_{dis,i,t}$ are introduced. The operational constraints for the BESS are defined as:
\begin{subequations}
\begin{align}
    0 \le P_{ch,i,t} &\le u_{ch,i,t} P_{BESS}^{\text{max}}, \quad \forall i \in \mathcal{N}^{BESS}, \forall t \label{eq:bess_char_limit} \\
    0 \le P_{dis,i,t} &\le u_{dis,i,t} P_{BESS}^{\text{max}}, \quad \forall i \in \mathcal{N}^{BESS}, \forall t \label{eq:bess_dis_limit} \\
    u_{ch,i,t} + u_{dis,i,t} &\le 1, \quad \forall i \in \mathcal{N}^{BESS}, \forall t \label{eq:bess_binary_limit}
\end{align}
\end{subequations}
where $P_{ch,i,t}$ and $P_{dis,i,t}$ are the charging and discharging power, respectively, bounded by the maximum capacity $P_{BESS}^{\text{max}}$.

The state-of-charge ($SOC$) of the battery updates dynamically across consecutive time intervals, constrained by the energy capacity ($E_{BESS}$) and the round-trip efficiency ($\eta_{BESS}$):
\begin{subequations}
\begin{align}
    SOC_{i,t} &= SOC_{i,t-1} + \frac{P_{ch,i,t}\eta_{BESS} - P_{dis,i,t}/\eta_{BESS}}{E_{BESS}}, \quad \forall t > 0 \label{eq:soc_update} \\
    SOC^{\text{min}} &\le SOC_{i,t} \le SOC^{\text{max}}, \quad \forall t \label{eq:soc_limit} \\
    SOC_{i,0} &= SOC_{i,24} = SOC^{\text{ini}} \label{eq:soc_equal}
\end{align}
\end{subequations}
Equation \eqref{eq:soc_equal} ensures the energy neutrality of the BESS over a 24-hour daily cycle, setting the initial and final $SOC$ to a predefined level $SOC^{\text{ini}}$.

\subsubsection{DG}
DGs operate as controllable power sources, governed by their minimum and maximum output boundaries. The operational status of the DG is indicated by a binary variable $u_{DG,i,t}$. To ensure the technical limits of the Diesel generators, additional constraints including reactive power limits, startup status, minimum operating times, and spinning reserve are defined as follows:
\begin{subequations}
\begin{align}
    u_{DG,i,t} P_{DG}^{\text{min}} \le P_{DG,i,t} &\le u_{DG,i,t} P_{DG}^{\text{max}}, \quad \forall i \in \mathcal{N}^{DG}, \forall t \label{eq:dg_active_limit} \\
    (P_{DG,i,t})^2 + (Q_{DG,i,t})^2 &\le (S_{DG}^{\text{max}})^2 u_{DG,i,t}, \quad \forall i \in \mathcal{N}^{DG}, \forall t \label{eq:dg_apparent_limit} \\
    u_{DG,i,t} Q_{DG}^{\text{min}} \le Q_{DG,i,t} &\le u_{DG,i,t} Q_{DG}^{\text{max}}, \quad \forall i \in \mathcal{N}^{DG}, \forall t \label{eq:dg_q_limit} \\
    v_{DG,i,t} &\ge u_{DG,i,t} - u_{DG,i,t-1}, \quad \forall i \in \mathcal{N}^{DG}, \forall t \label{eq:dg_startup} \\
    \sum_{\tau=t}^{t+T_{up}^{DG}-1} u_{DG,i,\tau} &\ge T_{up}^{DG} v_{DG,i,t}, \quad \forall i \in \mathcal{N}^{DG}, \forall t \label{eq:dg_uptime} \\
    \sum_{\tau=t}^{t+T_{down}^{DG}-1} (1 - u_{DG,i,\tau}) &\ge T_{down}^{DG} (u_{DG,i,t-1} - u_{DG,i,t}), \quad \forall i \in \mathcal{N}^{DG}, \forall t \label{eq:dg_downtime} \\
    R_{DG,t} &= \sum_{i \in \mathcal{N}^{DG}} \left( u_{DG,i,t} P_{DG}^{\text{max}} - P_{DG,i,t} \right), \quad \forall t \label{eq:dg_reserve_def} \\
    R_{DG,t} &\ge 0.2 \sum_{i \in \mathcal{N}^{DG}} P_{DG}^{\text{max}}, \quad \forall t \label{eq:dg_reserve_req}
\end{align}
\end{subequations}
where $P_{DG,i,t}$ and $Q_{DG,i,t}$ are the active and reactive power of the DG, respectively, and $S_{DG}^{\text{max}}$ represents its maximum apparent power. $v_{DG,i,t}$ denotes the startup state, and $R_{DG,t}$ represents the spinning reserve. The minimum up and down times are defined by $T_{up}^{DG}$ and $T_{down}^{DG}$.

Additionally, the variation in the DG's power output between consecutive time intervals is bounded by its ramp-up and ramp-down limits to prevent mechanical stress:
\begin{equation}
    -R_{DG} \le P_{DG,i,t} - P_{DG,i,t-1} \le R_{DG}, \quad \forall t > 0 \label{eq:dg_ramp_limit}
\end{equation}
where $R_{DG}$ is the maximum allowable ramping rate.

\subsection{AC Optimal Power Flow Formulation}
\label{subsec:ac_opf}

AC-OPF model. Operating at the foundational level of the hierarchy, the local AC-OPF problem is formulated to robustly handle both standard operation and unexpected fault scenarios (such as generator failures, extreme renewable drops, or main grid blackouts). To manage this, the system utilizes a binary emergency indicator $\Gamma_E(t) \in \{0, 1\}$ to seamlessly toggle between normal operation ($\Gamma_E(t) = 0$) and an emergency fault state ($\Gamma_E(t) = 1$), ensuring operational continuity and optimal economic dispatch under any physical condition.

\subsubsection{Economic and Loss Functions}
\label{subsubsec:economic_loss_functions}

The primary objective of the local microgrid operator is to minimize the total operational cost over the scheduling horizon $\mathcal{T}$. To generalize the optimization problem across varying grid conditions, the fundamental objective function unifies standard economic costs and load shedding penalty costs.

The total operational cost function $\mathcal{F}^{OPF}$ is expressed as:
\begin{equation}
    \mathcal{F}^{OPF} = F_{grid} + F_{DG} + F_{BESS} + F_{loss} + F_{shed}
    \label{eq:total_cost_function}
\end{equation}

The individual economic and operational cost components are detailed as follows:

\begin{subequations}
\label{eq:cost_components}
\begin{align}
    F_{grid} &= \sum_{t \in \mathcal{T}} \left( \pi_{buy,t} P_{grid,t}^{buy} + \pi_{sell,t} P_{grid,t}^{sell} \right) \label{eq:cost_grid} \\
    F_{DG} &= \sum_{t \in \mathcal{T}} \sum_{i \in \mathcal{N}^{DG}} \left( \alpha_{DG} u_{DG,i,t} + (\beta_{DG} + \kappa_{DG}) P_{DG,i,t} + \gamma_{DG} (P_{DG,i,t})^2 + C_{start,DG} v_{DG,i,t} \right) \label{eq:cost_dg} \\
    F_{BESS} &= \sum_{t \in \mathcal{T}} \sum_{i \in \mathcal{N}^{BESS}} C_{BESS} \left( \eta_{BESS} P_{ch,i,t} + \frac{P_{dis,i,t}}{\eta_{BESS}} \right) \label{eq:cost_bess} \\
    F_{loss} &= \sum_{t \in \mathcal{T}} \sum_{(i,j) \in \mathcal{E}} \omega_{loss} I_{ij,t}^2 \label{eq:cost_loss} \\
    F_{shed} &= \sum_{t \in \mathcal{T}} \sum_{i \in \mathcal{N}} \left( C_{shed}^{normal} P_{shed,i,t}^{normal} + C_{shed}^{critical} P_{shed,i,t}^{critical} \right) \label{eq:cost_shed}
\end{align}
\end{subequations}

Here, $F_{grid}$ represents the cost of purchasing electricity and the revenue from selling power to the main grid. $F_{DG}$ represents the operational cost of Diesel Generators, encompassing fixed, linear, and quadratic fuel costs, linear Operation and Maintenance (O\&M) costs (represented by coefficient $\kappa_{DG}$), and startup costs. $F_{BESS}$ calculates the degradation cost of the battery storage system proportional to its charging and discharging energy. To promote localized consumption and minimize distribution losses, $F_{loss}$ represents the active power losses along the distribution lines, weighted by a factor $\omega_{loss}$. 

Finally, $F_{shed}$ leverages static, prohibitively high VOLL multipliers ($C_{shed}^{normal}$, $C_{shed}^{critical}$) to strictly penalize load curtailment and prioritize critical load preservation. 

\subsubsection{Operational Modes and Fault Scenarios}
\label{subsubsec:operational_modes}

Rather than embedding fault logic directly into the physical equations (such as the DistFlow boundaries), the system isolates unexpected events as explicit operational constraints. Depending on the operational state, denoted by $\Gamma_E(t)$ (introduced in Section \ref{subsec:ac_opf}), specific mathematical conditions are enforced to reflect either standard operation or component failures.

To protect against severe power deficits during faults, the load shedding limits are strictly defined by their physical capacity. Additionally, to maintain a constant power factor, the reactive power shed is constrained to be strictly proportional to the active power shed:
\begin{subequations}
\begin{align}
    P_{shed,j,t} &= P_{shed,j,t}^{normal} + P_{shed,j,t}^{critical} \label{eq:shed_total} \\
    0 \le P_{shed,j,t}^{normal} &\le 0.8 P_{load,j,t} \label{eq:shed_normal_limit} \\
    0 \le P_{shed,j,t}^{critical} &\le 0.2 P_{load,j,t} \label{eq:shed_critical_limit} \\
    Q_{shed,j,t} &= P_{shed,j,t} \frac{Q_{load,j,t}}{P_{load,j,t}} \label{eq:shed_q_ratio}
\end{align}
\end{subequations}

The transition between modes is mathematically handled as follows:

\begin{itemize}
    \item \textbf{Normal Mode ($\Gamma_E(t) = 0$):} The microgrid operates with all components fully available. Due to the prohibitive VOLL penalties ($C_{shed}^{critical} \gg C_{shed}^{normal} \gg \text{cost}_{buy}$) in the objective function \eqref{eq:cost_shed}, the shedding variables are naturally forced to zero by the optimizer. The system maintains its standard security margins, including the required spinning reserve as defined in \eqref{eq:dg_reserve_req}.

    \item \textbf{Emergency Mode ($\Gamma_E(t) = 1$):} The system experiences a severe disruption, which could encompass a main grid blackout, a local generator trip, or an extreme drop in renewable generation. The failed component $k$ is mathematically isolated by forcing its power output and exchange variables to zero:
    \begin{equation}
        P_{k,t} = 0, \quad Q_{k,t} = 0 \quad \forall k \in \Omega_{fault} \label{eq:fault_generic}
    \end{equation}
    where $\Omega_{fault}$ represents the set of components experiencing an outage (e.g., the main grid interconnection, specific DGs, or PV units). 
    
    To counteract the resulting power deficit, the standard spinning reserve constraint \eqref{eq:dg_reserve_req} is intentionally relaxed. This relaxation allows the remaining dispatchable generators to fully utilize their 20\% locked reserve capacity to maintain grid stability. If the internal generation remains insufficient, the solver is permitted to shed load according to \eqref{eq:shed_normal_limit}-\eqref{eq:shed_critical_limit}, strongly guided by the economic prioritization in $F_{shed}$ to preserve critical loads.
\end{itemize}


\subsubsection{Distribution Power Flow Constraints}
\label{subsubsec:distflow}

The power flow within the radial network of the microgrids is modeled using the original non-linear Branch Flow Model (DistFlow) equations. The active and reactive power balance, voltage drop, and branch current definitions are formulated as follows:

\begin{subequations} \label{eq:distflow_original}
\begin{align}
    & \begin{aligned} \label{eq:P_balance}
        P_{line,ij,t} - R_{ij} I_{ij,t}^2 &+ P_{DG,j,t} + P_{PV,j,t} + P_{WT,j,t} + P_{dis,j,t} - P_{ch,j,t} \\
        &+ P_{tie,j,t} + P_{shed,j,t} - P_{load,j,t} = \sum_{k \in \Omega(j)} P_{line,jk,t}
    \end{aligned} \quad \forall j \in \mathcal{N}, \forall t \in \mathcal{T} \\
    & \begin{aligned} \label{eq:Q_balance}
        Q_{line,ij,t} - X_{ij} I_{ij,t}^2 &+ Q_{DG,j,t} \\
        &+ Q_{shed,j,t} - Q_{load,j,t} = \sum_{k \in \Omega(j)} Q_{line,jk,t}
    \end{aligned} \quad \forall j \in \mathcal{N}, \forall t \in \mathcal{T} \\
    & V_{j,t}^2 = V_{i,t}^2 - 2(R_{ij}P_{line,ij,t} + X_{ij}Q_{line,ij,t}) + (R_{ij}^2 + X_{ij}^2)I_{ij,t}^2 \quad \forall (i,j) \in \mathcal{E}, \forall t \in \mathcal{T} \label{eq:voltage_drop} \\
    & I_{ij,t}^2 = \frac{P_{line,ij,t}^2 + Q_{line,ij,t}^2}{V_{i,t}^2} \quad \forall (i,j) \in \mathcal{E}, \forall t \in \mathcal{T} \label{eq:current_def} \\
    & (V^{\text{min}})^2 \leq V_{j,t}^2 \leq (V^{\text{max}})^2 \quad \forall j \in \mathcal{N}, \forall t \in \mathcal{T} \label{eq:voltage_limit} \\
    & 0 \leq I_{ij,t}^2 \leq (I^{\text{max}})^2 \quad \forall (i,j) \in \mathcal{E}, \forall t \in \mathcal{T} \label{eq:current_limit}
\end{align}
\end{subequations}

To correctly model the interaction with the main grid, a boundary condition is enforced at the slack bus (denoted as node $j=0$). The branch power injected into the root node from the external network explicitly represents the grid power exchange:
\begin{subequations} \label{eq:distflow_boundary}
\begin{align}
    P_{line,0,t} &= P_{grid,t}^{buy} + P_{grid,t}^{sell} \quad \forall t \in \mathcal{T} \label{eq:boundary_P} \\
    Q_{line,0,t} &= Q_{grid,t} \quad \forall t \in \mathcal{T} \label{eq:boundary_Q}
\end{align}
\end{subequations}

Equations \eqref{eq:P_balance} and \eqref{eq:Q_balance} delineate the active and reactive power balance at node $j$ at time $t$, respectively. The sum of all active power injected into node $j$---including the incoming branch power ($P_{line,ij,t}$), distributed generation ($P_{DG,j,t}$), renewable active power generation ($P_{PV,j,t}, P_{WT,j,t}$), BESS net discharge ($P_{dis,j,t} - P_{ch,j,t}$), P2P trading ($P_{tie,j,t}$), and load shedding ($P_{shed,j,t}$)---must equal the sum of the power flowing out to child branches $\sum P_{line,jk,t}$, the local load demand ($P_{load,j,t}$), and the line active power losses ($R_{ij} I_{ij,t}^2$). By convention, active power sold to the grid ($P_{grid,t}^{sell}$) takes non-positive values. 

Similarly, the reactive power balance in \eqref{eq:Q_balance} accounts for branch flows ($Q_{line,ij,t}$), DG reactive support ($Q_{DG,j,t}$), and reactive load shedding ($Q_{shed,j,t}$). Notably, the renewable energy inverters (PV and WT) are configured to operate at unity power factor, contributing zero reactive power to the system.

Equation \eqref{eq:voltage_drop} governs the voltage drop along branch $(i,j)$, considering the line impedance ($R_{ij}, X_{ij}$) and branch flows. Equation \eqref{eq:current_def} represents the non-linear definition of the squared branch current. Finally, inequalities \eqref{eq:voltage_limit} and \eqref{eq:current_limit} ensure that the squared nodal voltages and branch currents remain within their secure operational limits. This exact DistFlow model is specifically chosen as it inherently captures the physical radial topology of the microgrids, forming the strict prerequisite for subsequent global optimization. However, due to the non-linear and non-convex nature of the current definition in \eqref{eq:current_def}, this problem requires convexification via SOCP to be mathematically tractable for optimization solvers.

\subsection{SOCP Relaxation and Exactness}
\label{subsec:socp_relaxation}

\subsubsection{Convexification and SOC Relaxation}
The non-linear nature of the AC Optimal Power Flow fundamentally stems from the current magnitude definition in the Branch Flow Model. As originally stated in Equation \eqref{eq:current_def}, the squared branch current is formulated as:
\begin{equation}
    I_{ij,t}^2 = \frac{P_{line,ij,t}^2 + Q_{line,ij,t}^2}{V_{i,t}^2} \nonumber
\end{equation}
This equality constraint is highly non-convex. To achieve a computationally tractable convex model, the SOC relaxation transforms the non-convex feasible region into a convex cone by replacing the equality with an inequality:
\begin{equation}
    I_{ij,t}^2 V_{i,t}^2 \ge P_{line,ij,t}^2 + Q_{line,ij,t}^2
    \label{eq:socp_relaxation}
\end{equation}

\subsubsection{Exactness in Normal Mode}
The fundamental premise of substituting an equality with an inequality is predicated on the concept of exactness---the condition where the optimizer natively forces the inequality \eqref{eq:socp_relaxation} to hold with equality. In the Normal Mode of operation, this exactness is inherently guaranteed. As defined in the objective function \eqref{eq:total_cost_function}, the formulation incorporates a dedicated loss penalty component, $F_{loss}$ \eqref{eq:cost_loss}. Because the primary objective of the solver is to minimize the total operational cost, it actively suppresses the squared current variable $I_{ij,t}^2$, strictly driving it towards its theoretical lower bound. Consequently, the inequality holds with equality ($I_{ij,t}^2 V_{i,t}^2 = P_{line,ij,t}^2 + Q_{line,ij,t}^2$), thereby guaranteeing the exactness of the SOCP model and accurately reflecting the true physical state of the distribution grid.

\subsubsection{Robustness against Inexactness in Emergency Mode}
While exactness is guaranteed under standard operations, the mathematical literature suggests that significant challenges can arise during extreme Emergency Modes. A typical vulnerability occurs during islanded operations with a significant surplus of renewable energy generation. In standard models lacking curtailment mechanisms, if the microgrid cannot export this excess power, the system risks violating the exactness ($I_{ij,t}^2 V_{i,t}^2 > P_{line,ij,t}^2 + Q_{line,ij,t}^2$), which renders the optimization results physically meaningless. This breakdown occurs because the mathematical solver might erroneously exploit the SOCP relaxation gap, artificially inflating the current variable $I_{ij,t}^2$ to burn the surplus energy as fictitious line losses.

However, the proposed microgrid formulation is intrinsically designed to effectively counteract this phenomenon. The renewable generation constraints natively allow for free active power curtailment \eqref{eq:pv_limit}. Confronted with excess generation, the solver will preferentially curtail the renewable output at zero cost rather than inflating $I_{ij,t}^2$, which would incur a strict financial penalty via the $F_{loss}$ term scaled by $\omega_{loss}$. This dual mechanism---free curtailment capability combined with a strict virtual loss penalty---ensures that the solver is mathematically prevented from exploiting the relaxation gap. As a result, the model robustly preserves the physical integrity and exactness of the network's power flow equations even amidst extreme emergency and over-generation scenarios.

\subsection{Final Local SOCP Optimization Model}
\label{subsec:final_local_socp}

Integrating the operational cost components from Section \ref{subsubsec:economic_loss_functions} and the relaxed physical network constraints, the final SOCP model is formulated. Although the problem is implemented and solved as a single unified code using the emergency indicator $\Gamma_E(t)$, from a mathematical presentation perspective, the active objective function and constraints can be categorized into two distinct formulations based on the operational state.

\subsubsection{Formulation I: Normal Operation Mode ($\Gamma_E(t) = 0$)}
In standard grid-connected mode, the load shedding penalty enforces zero curtailment ($F_{shed} \to 0$). The objective is purely focused on economic costs and loss minimization. The problem is defined as:
\begin{equation}
\begin{aligned}
    \min \quad & \mathcal{F}_{normal}^{OPF} = F_{grid} + F_{DG} + F_{BESS} + F_{loss} \\
    \text{s.t.} \quad & \text{Component Constraints } \eqref{eq:pv_limit}-\eqref{eq:dg_ramp_limit} \\
    & \text{Standard Spinning Reserve } \eqref{eq:dg_reserve_req} \\
    & \text{DistFlow Boundary Conditions } \eqref{eq:distflow_boundary} \\
    & \text{Power Balance \& Voltage Constraints } \eqref{eq:P_balance}-\eqref{eq:voltage_drop}, \eqref{eq:voltage_limit}-\eqref{eq:current_limit} \\
    & \text{SOCP Relaxed Current Constraint } \eqref{eq:socp_relaxation}
\end{aligned}
\label{eq:total_objective_normal}
\end{equation}

\subsubsection{Formulation II: Emergency Mode ($\Gamma_E(t) = 1$)}
When a generalized fault occurs, power generation or exchange with the affected components $k \in \Omega_{fault}$ completely ceases. Consequently, the operational costs associated with these specific failed components are eliminated from the objective function (as their power outputs are forced to zero). The primary goal of the remaining microgrid transitions from pure economic profit to system resilience, leveraging the unlocked reserve capacity and prioritizing critical load preservation. The problem transitions to:
\begin{equation}
\begin{aligned}
    \min \quad & \mathcal{F}_{emergency}^{OPF} = F_{grid} + F_{DG} + F_{BESS} + F_{loss} + F_{shed} \\
    \text{s.t.} \quad & \text{Component Constraints } \eqref{eq:pv_limit}-\eqref{eq:dg_ramp_limit} \\
    & \text{Load Shedding Limits } \eqref{eq:shed_total}-\eqref{eq:shed_q_ratio} \\
    & \text{Operational Fault Constraints } \eqref{eq:fault_generic} \\
    & \text{DistFlow Boundary Conditions } \eqref{eq:distflow_boundary} \\
    & \text{Power Balance \& Voltage Constraints } \eqref{eq:P_balance}-\eqref{eq:voltage_drop}, \eqref{eq:voltage_limit}-\eqref{eq:current_limit} \\
    & \text{SOCP Relaxed Current Constraint } \eqref{eq:socp_relaxation}
\end{aligned}
\label{eq:total_objective_emergency}
\end{equation}

